\documentclass[11pt,a4paper]{article}

\usepackage[T1]{fontenc}
\usepackage[utf8]{inputenc}
\usepackage{lmodern}
\usepackage[french]{babel}
\frenchsetup{SmallCapsFigTabCaptions=false}
\addto\captionsfrench{\renewcommand{\tablename}{Tableau}}
\usepackage[a4paper,margin=2.5cm]{geometry}
\usepackage{amsmath}
\usepackage{graphicx}
\usepackage{float}
\usepackage{booktabs}
\usepackage[locale=FR]{siunitx}
\usepackage{caption}
\usepackage{xcolor}
\usepackage{tikz}
\definecolor{bleuENP}{RGB}{34,64,154}
\usepackage{listings}
\usepackage{hyperref}

\captionsetup{font=small,labelfont=bf}
\hypersetup{colorlinks=true,linkcolor=black,urlcolor=black}

\lstset{language=Matlab, basicstyle=\ttfamily\scriptsize, numbers=left,
  numberstyle=\tiny, numbersep=6pt, frame=single, breaklines=true,
  columns=fullflexible, keepspaces=true, showstringspaces=false,
  commentstyle=\color{gray}}

% Figure : \fig{fichier sans extension}{légende}{largeur en fraction de \linewidth}
\newcommand{\fig}[3]{\begin{figure}[H]\centering\includegraphics[width=#3\linewidth]{figures/#1}\caption{#2}\label{fig:#1}\end{figure}}
% Exemple pour une photo ou une capture d'écran (fichier figures/photo_banc.jpg) :
% \fig{photo_banc.jpg}{Montage en salle de TP.}{0.6}
% puis dans le texte : voir la figure~\ref{fig:photo_banc.jpg}.

% Deux figures côte à côte : \figdeux{fichier 1}{légende 1}{fichier 2}{légende 2}
\newcommand{\figdeux}[4]{\begin{figure}[H]\centering
\begin{minipage}[t]{0.49\linewidth}\centering\includegraphics[width=0.92\linewidth]{figures/#1}\caption{#2}\label{fig:#1}\end{minipage}\hfill
\begin{minipage}[t]{0.49\linewidth}\centering\includegraphics[width=0.92\linewidth]{figures/#3}\caption{#4}\label{fig:#3}\end{minipage}\end{figure}}

% Valeurs imprimées par TP01_ligne_longue.m (exécution du script).
% Si une valeur change après une nouvelle exécution, la modifier ici seulement.

% --- Constantes de la ligne ---
\newcommand{\Zc}{273{,}86}
\newcommand{\vp}{182\,574}
\newcommand{\tauL}{547{,}72}

% --- Contrôle par rapport au livre (tau = 0,00055 s) ---
\newcommand{\iunA}{685{,}2}
\newcommand{\vunA}{187\,661}
\newcommand{\ideuxA}{0{,}0}
\newcommand{\vdeuxA}{0}
\newcommand{\iunB}{683{,}8}
\newcommand{\vunB}{187\,261}
\newcommand{\ideuxB}{0{,}0}
\newcommand{\vdeuxB}{0}
\newcommand{\iunC}{662{,}0}
\newcommand{\vunC}{181\,293}
\newcommand{\ideuxC}{66{,}1}
\newcommand{\vdeuxC}{356\,825}

% --- Q1 : régime permanent (simulation / phaseurs) ---
\newcommand{\IunSim}{429{,}0}
\newcommand{\IunPh}{428{,}9}
\newcommand{\IdeuxSim}{448{,}5}
\newcommand{\IdeuxPh}{448{,}5}
\newcommand{\VdeuxSim}{184{,}3}
\newcommand{\VdeuxPh}{184{,}3}

% --- Q2, Q3 : coefficients de réflexion et rapports max|v2|/Vm ---
\newcommand{\GRL}{0{,}1872}
\newcommand{\Gneuf}{0{,}5697}
\newcommand{\Gmega}{0{,}9995}
\newcommand{\Gcc}{-1{,}0000}
\newcommand{\rapNeuf}{1{,}57}
\newcommand{\rapMega}{2{,}01}
\newcommand{\rapCC}{0{,}136}

% --- Q5 : source continue, valeurs à t = 5 s ---
\newcommand{\dcVRL}{230{,}0}
\newcommand{\dcVOuv}{230{,}0}
\newcommand{\dcVCC}{8{,}96}
\newcommand{\dcIRL}{575{,}0}
\newcommand{\dcIOuv}{230{,}2}
\newcommand{\dcICC}{7\,491\,036}

% --- Q6 : onde de choc ---
\newcommand{\chocCrete}{206{,}7}
\newcommand{\chocTc}{3{,}99}
\newcommand{\etaVal}{1{,}1125}
\newcommand{\fVmax}{0}
\newcommand{\fImax}{0}
\newcommand{\gArrRL}{54{,}78}
\newcommand{\gArrOuv}{54{,}78}
\newcommand{\gArrCC}{54{,}78}
\newcommand{\gVRL}{458{,}6}
\newcommand{\gVOuv}{369{,}2}
\newcommand{\gVCC}{0{,}0}

% --- Valeurs calculées à la main à partir des valeurs ci-dessus et des données ---
\newcommand{\deuxTau}{1{,}095}         % 2 tau en ms
\newcommand{\marcheDC}{1\,680}          % 2 x 230 kV / Zc en A
\newcommand{\nbAR}{4\,564}              % nombre d'allers-retours en 5 s
\newcommand{\estimDC}{7{,}7}            % 4564 x 1680 A en MA
\newcommand{\GneufV}{361}               % (1 + Gamma(999)) x 230 kV en kV
\newcommand{\ReqChoc}{33}               % 2 x 0,001 H / 60 ns en kilo-ohms
\newcommand{\ReqCC}{20}                 % 2 x 0,001 H / 0,1 ms en ohms

% --- Valeurs du livre (solution.pdf) ---
\newcommand{\livreiunA}{685}
\newcommand{\livrevunA}{187\,661}
\newcommand{\livreiunB}{683}
\newcommand{\livrevunB}{187\,261}
\newcommand{\livreiunC}{662}
\newcommand{\livrevunC}{181\,293}
\newcommand{\livreideuxC}{66}
\newcommand{\livrevdeuxC}{356\,731}


\begin{document}

% Page de garde : tout est placé en coordonnées absolues (cm) depuis le coin bas gauche.
\begin{titlepage}
\thispagestyle{empty}
\begin{tikzpicture}[remember picture,overlay,x=1cm,y=1cm]
\begin{scope}[shift={(current page.south west)}]

% --- bloc bleu supérieur, coupé en diagonale ---
\shade[top color=bleuENP!70!black, bottom color=bleuENP]
  (0,29.7) -- (21,29.7) -- (21,15.2) -- (0,12.6) -- cycle;
\fill[orangeTP] (0,12.6) -- (21,15.2) -- (21,14.95) -- (0,12.35) -- cycle;

% --- champ d'ondes progressives en filigrane ---
\begin{scope}
  \clip (0,29.7) -- (21,29.7) -- (21,15.2) -- (0,12.6) -- cycle;
  \foreach \k in {0,...,13} {
    \draw[white, opacity=0.05+0.012*\k, line width=0.5pt, smooth, domain=0:21, samples=120]
      plot (\x, {14.2+1.05*\k + (0.25+0.03*\k)*sin(\x*(38+4*\k) + 37*\k)});
  }
\end{scope}

% --- grand « 01 » en transparence ---
\node[anchor=north east, white, opacity=0.10, inner sep=0, font=\sffamily\bfseries\fontsize{170}{170}\selectfont]
  at (20.4,29.2) {01};

% --- logo sur carte blanche ---
\fill[white, rounded corners=6pt] (1.5,24.6) rectangle (6.2,28.7);
\node at (3.85,26.65) {\includegraphics[width=4.0cm]{figures/logo_enp.png}};

% --- surtitre et titre ---
\node[anchor=west, orangeTP, font=\sffamily\bfseries\small] at (1.5,23.2)
  {\textls[200]{TRAVAUX PRATIQUES \,·\, RÉSEAUX ÉLECTRIQUES}};
\node[anchor=west, white, font=\sffamily\bfseries\fontsize{50}{50}\selectfont] at (1.35,21.5)
  {Propagation};
\node[anchor=west, white, font=\sffamily\fontsize{25}{28}\selectfont] at (1.5,19.75)
  {dans les lignes électriques};
\fill[orangeTP] (1.55,18.75) rectangle (4.6,18.9);

% --- signal : réponse en escalier de v2 à l'extrémité ouverte ---
\draw[orangeTP!40, line width=5pt, opacity=0.35, line join=round]
  (1.5,15.25) -- (4.0,15.25) -- (4.0,17.35) -- (6.2,17.35) -- (6.2,16.1) -- (8.4,16.1)
  -- (8.4,16.95) -- (10.6,16.95) -- (10.6,16.4) -- (12.8,16.4) -- (12.8,16.75) -- (15.0,16.75)
  -- (15.0,16.52) -- (17.2,16.52) -- (17.2,16.66) -- (19.5,16.66);
\draw[orangeTP, line width=1.6pt, line join=round]
  (1.5,15.25) -- (4.0,15.25) -- (4.0,17.35) -- (6.2,17.35) -- (6.2,16.1) -- (8.4,16.1)
  -- (8.4,16.95) -- (10.6,16.95) -- (10.6,16.4) -- (12.8,16.4) -- (12.8,16.75) -- (15.0,16.75)
  -- (15.0,16.52) -- (17.2,16.52) -- (17.2,16.66) -- (19.5,16.66);
\fill[white] (4.0,17.35) circle (2.2pt);
\node[anchor=south west, white, font=\sffamily\footnotesize, opacity=0.85] at (4.1,17.45)
  {$v_2 \approx (1+\Gamma)\,v_{\mathrm{inc}}$};
\node[anchor=north, white, font=\sffamily\scriptsize, opacity=0.8] at (5.1,16.0) {$2\tau$};
\draw[white, opacity=0.8, |<->|, line width=0.5pt] (4.0,16.3) -- (6.2,16.3);

% --- sous-titre ---
\node[anchor=west, text width=17.5cm, bleuENP!70!black, font=\large\itshape] at (1.5,10.2)
  {Modèle de Bergeron et règle des trapèzes : régimes transitoires\\ d'une ligne longue sans pertes};

% --- cartes d'information ---
\foreach \x/\titre/\ligneA/\ligneB in {
    1.5/RÉALISÉ PAR/Mohammed ZANE/Groupe 1,
    8.0/ENCADRÉ PAR/Prof. Karim SEBAA/Réseaux Électriques,
    14.5/ANNÉE/2026 -- 2027/TP n\textsuperscript{o}\,01} {
  \fill[bleuENP!6] (\x,5.5) rectangle ++(5.2,2.9);
  \fill[orangeTP] (\x,5.5) rectangle ++(0.14,2.9);
  \node[anchor=north west, bleuENP, font=\sffamily\bfseries\scriptsize] at (\x+0.45,8.05)
    {\textls[180]{\titre}};
  \node[anchor=north west, black, font=\sffamily\bfseries\large] at (\x+0.45,7.35) {\ligneA};
  \node[anchor=north west, black!65, font=\sffamily\small] at (\x+0.45,6.5) {\ligneB};
}

% --- pied de page ---
\draw[bleuENP!25, line width=0.4pt] (1.5,3.1) -- (19.5,3.1);
\node[anchor=west, black!60, font=\sffamily\small] at (1.5,2.5)
  {École Nationale Polytechnique \,·\, Département de Génie Électrique};
\node[anchor=east, black!60, font=\sffamily\small] at (19.5,2.5) {Alger};
\fill[bleuENP] (0,0) rectangle (21,0.55);
\fill[orangeTP] (0,0.55) rectangle (7,0.7);

\end{scope}
\end{tikzpicture}
\null
\end{titlepage}


\section{Objectif}

On simule la mise sous tension d'une ligne monophasée sans pertes de longueur
$d = 100$~mi ($L' = 1{,}5$~mH/mi, $C' = 0{,}02$~\unit{\micro\farad}/mi) alimentée par
$v_s(t) = 230\,000\,\sqrt{2}/\sqrt{3}\,\cos(2\pi\,60\,t)$ et fermée sur une charge R-L série
($R = \SI{400}{\ohm}$, $L = \SI{0.25}{H}$). L'interrupteur se ferme à $t_f = \SI{0.1}{ms}$ et
le pas de calcul vaut $\Delta t = \SI{0.1}{ms}$, avec des conditions initiales nulles.
On reprend ensuite le calcul pour une extrémité ouverte, une extrémité en court-circuit, une
source continue de 230~kV et une onde de choc de foudre. Les grandeurs suivies sont $i_1$, $v_1$
au départ (extrémité $k$) et $i_2$, $v_2$ à l'arrivée (extrémité $m$).

\fig{circuit}{Circuit étudié : source, interrupteur fermé à \SI{0.1}{ms}, ligne de 100~mi et
charge R-L.}{0.62}

\section{Méthode}

% --- Modèle de Bergeron ---
Pour la ligne sans pertes, l'impédance caractéristique, la vitesse de propagation et le temps de
propagation valent
\begin{equation}
Z_c = \sqrt{L'/C'}, \qquad v_p = 1/\sqrt{L'C'}, \qquad \tau = d/v_p ,
\end{equation}
soit $Z_c = \Zc~\unit{\ohm}$, $v_p = \vp$~mi/s et $\tau = \tauL~\unit{\micro s}$.
Le modèle de Bergeron remplace chaque extrémité par la résistance $Z_c$ en parallèle avec une
source de courant d'historique. Cette source ne dépend que de l'autre extrémité à l'instant
$t-\tau$ ; elle est donc connue quand on calcule l'instant $t$. Avec $i_1$ entrant dans la ligne
en $k$ et $i_2$ sortant en $m$ vers la charge :
\begin{align}
i_1(t) &= \frac{v_1(t)}{Z_c} + I_k, & I_k &= i_2(t-\tau) - \frac{v_2(t-\tau)}{Z_c}, \label{eq:k}\\
i_2(t) &= I_m - \frac{v_2(t)}{Z_c}, & I_m &= i_1(t-\tau) + \frac{v_1(t-\tau)}{Z_c}. \label{eq:m}
\end{align}

% --- Modèle compagnon de l'inductance ---
L'inductance de charge est traitée par la règle des trapèzes. L'équation
$\mathrm{d}i_L/\mathrm{d}t = v_L/L$ devient un circuit linéaire, le modèle compagnon : une
résistance $R_{eq}$ en parallèle avec une source $I_h$ connue au pas précédent,
\begin{equation}
i_L(t) = \frac{v_L(t)}{R_{eq}} + I_h, \qquad R_{eq} = \frac{2L}{\Delta t}, \qquad
I_h = i_L(t-\Delta t) + \frac{v_L(t-\Delta t)}{R_{eq}} ,
\end{equation}
avec ici $R_{eq} = \SI{5000}{\ohm}$. Côté $m$, $Z_c \parallel I_m$ alimente $R$ en série avec
$R_{eq} \parallel I_h$. La résolution de ce petit circuit « continu » donne directement
\begin{equation}
i_2 = \frac{I_m + R_{eq} I_h/Z_c}{1 + (R+R_{eq})/Z_c}, \qquad v_2 = Z_c\,(I_m - i_2), \qquad
v_L = v_2 - R\,i_2 .
\end{equation}
Côté $k$, la tension est imposée : $v_1 = v_s$ dès que l'interrupteur est fermé, puis $i_1$ par
\eqref{eq:k}.

% --- Interpolation du retard ---
Le rapport $\tau/\Delta t$ n'est pas entier, donc $t-\tau$ tombe entre deux échantillons
stockés. La valeur retardée s'obtient par interpolation linéaire :
\begin{equation}
x(t-\tau) \approx x_j + a\,(x_{j+1}-x_j), \qquad j = \left\lfloor \frac{t-\tau}{\Delta t}
\right\rfloor, \qquad a = \frac{t-\tau}{\Delta t} - j .
\end{equation}
Il ne faut pas interpoler à travers l'instant de fermeture. Si $t-\tau < t_f$, la ligne était
encore au repos et l'historique vaut zéro. Sinon, l'interpolation entre l'échantillon à $t=0$
(interrupteur ouvert) et celui à $t_f$ ferait arriver l'onde en $m$ un pas trop tôt.

\section{Validation}

% --- Comparaison avec le livre ---
Le script refait d'abord les premiers pas de l'exemple du livre, avec $\tau$ arrondi à
\SI{0.55}{ms} comme dans la solution.

\begin{table}[H]
\centering
\caption{Charge R-L : valeurs du livre et du script ($\tau = \SI{0.55}{ms}$).}
\label{tab:livre}
\begin{tabular}{llrrrr}
\toprule
$t$ (ms) & source & $i_1$ (A) & $v_1$ (V) & $i_2$ (A) & $v_2$ (V) \\
\midrule
0,1 & livre  & \livreiunA & \livrevunA & 0 & 0 \\
    & script & \iunA & \vunA & \ideuxA & \vdeuxA \\
0,2 & livre  & \livreiunB & \livrevunB & 0 & 0 \\
    & script & \iunB & \vunB & \ideuxB & \vdeuxB \\
0,7 & livre  & \livreiunC & \livrevunC & \livreideuxC & \livrevdeuxC \\
    & script & \iunC & \vunC & \ideuxC & \vdeuxC \\
\bottomrule
\end{tabular}
\end{table}

Les valeurs coïncident à l'arrondi près. Le seul écart visible porte sur $v_2$ à
\SI{0.7}{ms} : \vdeuxC~V contre \livrevdeuxC~V. Le livre arrondit $Z_c$ à \SI{274}{\ohm} et
garde des valeurs intermédiaires arrondies (la source d'historique $I_m$ est écrite 1368~A),
alors que le script garde $Z_c = \Zc~\unit{\ohm}$ et toutes les décimales. L'onde arrive en $m$
au même pas que dans le livre, \SI{0.7}{ms}. La suite du calcul utilise le $\tau$ exact,
$\tauL~\unit{\micro s}$.

% --- Régime permanent par phaseurs ---
Le régime permanent de Q1 se calcule à la main en écrivant la ligne comme un quadripôle ABCD
sans pertes ($A = D = \cos\beta d$, $B = jZ_c\sin\beta d$, $C = j\sin\beta d/Z_c$, avec
$\beta d = \omega\tau$) fermé sur $Z_L = R + jL\omega$. Le tableau~\ref{tab:phaseurs} compare
ces amplitudes aux crêtes simulées entre 4,9 et 5~s.

\begin{table}[H]
\centering
\caption{Q1 : amplitudes en régime permanent, simulation et phaseurs.}
\label{tab:phaseurs}
\begin{tabular}{lrr}
\toprule
grandeur & simulation & phaseurs \\
\midrule
$|I_1|$ (A)  & \IunSim & \IunPh \\
$|I_2|$ (A)  & \IdeuxSim & \IdeuxPh \\
$|V_2|$ (kV) & \VdeuxSim & \VdeuxPh \\
\bottomrule
\end{tabular}
\end{table}

L'accord est au dixième près. Ce contrôle vérifie aussi le signe de $I_k$ : l'onde réfléchie ne
revient en $k$ qu'après $2\tau$, elle n'apparaît pas dans le tableau~\ref{tab:livre}, mais une
erreur de signe fausserait complètement le régime permanent.

\section{Résultats}

% --- Q1 : charge R-L ---
\subsection{Q1 : charge R-L}

À la fermeture, la source ne voit que l'impédance caractéristique : $i_1$ saute à \iunA~A,
soit $v_1/Z_c$. Rien ne bouge en $m$ tant que l'onde incidente n'est pas arrivée. Elle parvient
en $m$ $\tauL~\unit{\micro s}$ après la fermeture, et le premier pas où elle est vue est
\SI{0.7}{ms}. À cet instant $v_2$ vaut \vdeuxC~V, près du double de l'onde incidente :
au front, l'inductance se comporte presque en circuit ouvert ($R + R_{eq} = \SI{5400}{\ohm}$).
L'onde réfléchie revient ensuite en $k$ tous les $2\tau = \deuxTau$~ms, ce qui donne les
sauts de $i_1$ de la figure~\ref{fig:Q1_RL_zoom}. Après quelques dizaines de millisecondes, les
grandeurs sont sinusoïdales, avec les amplitudes du tableau~\ref{tab:phaseurs}.

\figdeux{Q1_RL_zoom}{Q1, charge R-L, 10 premières ms : échelon de $i_1$ à la fermeture, arrivée
de l'onde en $m$ à \SI{0.7}{ms}, sauts de $i_1$ à chaque retour de l'onde réfléchie.}
{Q1_RL_5s}{Q1, charge R-L, 0 à 5~s : régime sinusoïdal établi après le transitoire.}

% --- Q2 : extrémité ouverte ---
\subsection{Q2 : extrémité ouverte}

Le coefficient de réflexion $\Gamma = (Z_L - Z_c)/(Z_L + Z_c)$ vaut \Gneuf{} pour
$R = \SI{999}{\ohm}$ et \Gmega{} pour $R = \SI{e6}{\ohm}$ (\GRL{} pour la charge de Q1). La
tension transmise en $m$ est $(1+\Gamma)$ fois l'onde incidente. Le rapport $\max|v_2|/V_m$
vaut \rapNeuf{} pour 999~\unit{\ohm} et \rapMega{} pour \SI{e6}{\ohm} : seule la seconde valeur
correspond au doublement de tension d'une vraie extrémité ouverte. Avec \SI{e6}{\ohm}, le
courant $i_2$ reste nul et $v_2$ oscille autour de la sinusoïde.

\figdeux{Q2_ouvert_zoom}{Q2, extrémité ouverte, 20 premières ms : à l'arrivée, $v_2$ monte à
\rapNeuf{}\,$V_m$ pour 999~$\Omega$ et à \rapMega{}\,$V_m$ pour $10^6$~$\Omega$.}
{Q2_ouvert_5s}{Q2, extrémité ouverte, 0 à 5~s : l'oscillation libre disparaît alors que la
ligne est sans pertes (voir la discussion).}

% --- Q3 : court-circuit ---
\subsection{Q3 : extrémité en court-circuit}

Pour $R = L = 0{,}001$, $\Gamma = \Gcc$ : l'onde réfléchie annule la tension et double le
courant. Le rapport $\max|v_2|/V_m$ ne vaut que \rapCC. Sur la figure~\ref{fig:Q3_cc_zoom},
$v_2$ change de signe à chaque pas après l'arrivée de l'onde ; ce point est repris dans la
discussion. Le courant $i_2$ monte par paliers au fil des allers-retours.

\figdeux{Q3_cc_zoom}{Q3, court-circuit, 20 premières ms : $v_2$ reste faible et alterne de signe
d'un pas à l'autre, $i_2$ croît par paliers.}
{Q3_cc_5s}{Q3, court-circuit, 0 à 5~s : courants sinusoïdaux établis, $v_2$ reste faible.}

% --- Q5 : source continue ---
\subsection{Q5 : source continue de 230~kV}

Avec la source continue, l'onde incidente est un échelon constant et chaque réflexion ajoute un
palier : les réponses deviennent des escaliers de largeur $2\tau$
(figure~\ref{fig:Q5_DC_zoom}). Pour la charge R-L et pour 999~\unit{\ohm}, $v_2$ tend vers la
tension de la source : à $t = \SI{5}{s}$, $v_2 = \dcVRL$~kV et \dcVOuv~kV, et $i_2 = \dcIRL$~A
et \dcIOuv~A, c'est-à-dire $V/R$.

En court-circuit, $i_2$ augmente d'environ $2V/Z_c = \marcheDC$~A à chaque aller-retour de
l'onde ($2\tau = \deuxTau$~ms). Sur 5~s, cela fait \nbAR{} allers-retours, soit environ
\estimDC~MA. Le script donne \dcICC~A à $t = \SI{5}{s}$, un peu moins, parce que la tension aux
bornes de $R$ et $L$ ($v_2 = \dcVCC$~kV à 5~s) réduit légèrement chaque marche.

\figdeux{Q5_DC_zoom}{Q5, source continue, 10 premières ms : escaliers de largeur $2\tau$ ; en
court-circuit $i_1$ et $i_2$ montent d'une marche à chaque aller-retour.}
{Q5_DC_5s}{Q5, charges R-L et 999~$\Omega$, 0 à 5~s : $v_2$ tend vers 230~kV et le courant
vers $V/R$.}

\fig{Q5_DC_cc_5s}{Q5, court-circuit, 0 à 5~s : le courant croît presque linéairement jusqu'à
\dcICC~A à 5~s.}{0.47}

% --- Q6 : onde de choc de foudre ---
\subsection{Q6 : onde de choc de foudre}

La formule de l'énoncé, avec $\tau_1 = \SI{1.2}{\micro s}$, $\tau_2 = \SI{50}{\micro s}$ et
$n = 3$, culmine à \chocCrete~kV à \chocTc~\unit{\micro s}, et non à $V_{max} = 230$~kV. On
multiplie l'onde par $\eta = V_{max}/\text{crête} = \etaVal$ pour ramener la crête à 230~kV
(figure~\ref{fig:Q6a_onde}). Les calculs sur la ligne utilisent l'onde corrigée, avec
$\Delta t = \SI{60}{ns}$ et une fenêtre de \SI{250}{\micro s}.

\fig{Q6a_onde}{Q6a, onde de choc : formule de l'énoncé et onde corrigée par $\eta$.}{0.45}

Sur la ligne de 100~mi, $\tau = \tauL~\unit{\micro s}$ dépasse la fenêtre de
\SI{250}{\micro s} : l'onde n'atteint jamais l'extrémité $m$. Le script donne
$\max|v_2| = \fVmax$~V et $\max|i_2| = \fImax$~A, et $i_1$ reproduit $v_s/Z_c$ pendant toute la
simulation. Le résultat ne dépend donc pas de la charge.

Pour voir les réflexions dans la même fenêtre, on refait le calcul sur 10~mi : le temps de
propagation devient $\tau/10$ et l'onde arrive en $m$ à \gArrRL~\unit{\micro s} pour les trois
charges. La charge R-L se comporte en circuit ouvert au front, car $R_{eq} = 2L/\Delta t$ est
très grand à \SI{60}{ns} : $v_2$ monte à \gVRL~kV, près du double de 230~kV. Avec
999~\unit{\ohm}, $v_2$ atteint \gVOuv~kV, au-dessus de $(1+\Gamma)\cdot 230 \approx \GneufV$~kV :
l'inductance de 0,001~H en série vaut $R_{eq} = 2L/\Delta t \approx \ReqChoc$~k$\Omega$ à
$\Delta t = \SI{60}{ns}$, si bien que la charge paraît presque ouverte au front. Le court-circuit
résistif ($R = 0{,}001$, $L = 0$) annule $v_2$ (\gVCC~kV). L'onde réfléchie revient en $k$ à
$2\tau/10$, puis repart vers $m$ avec le signe opposé après réflexion sur la source.

\figdeux{Q6_100mi}{Q6, ligne de 100~mi : $i_1$ suit la forme de $v_s$, $v_2$ et $i_2$ restent
nuls pendant les \SI{250}{\micro s}.}
{Q6_10mi}{Q6, ligne de 10~mi : doublement de $v_2$ pour la charge R-L, réflexions successives à
intervalles de $2\tau/10$.}

\section{Discussion (Q4)}

% --- Propagation et réflexions ---
Dans tous les cas, l'extrémité $m$ ne réagit qu'après le temps de propagation $\tau$, et la
source ne voit d'abord que $Z_c$, quelle que soit la charge. C'est le coefficient de réflexion qui
fixe ensuite ce qui se passe en $m$ : la tension transmise vaut $(1+\Gamma)$ fois l'onde
incidente. Elle double presque à l'extrémité ouverte et s'annule en court-circuit, où c'est le
courant qui double. La source idéale impose $v_1$ ; pour l'onde réfléchie elle se comporte comme
un court-circuit ($\Gamma = -1$ en $k$), si bien que l'onde change de signe à chaque retour en $k$.

% --- 999 ohms ---
La résistance de 999~\unit{\ohm} proposée dans l'énoncé n'est pas une vraie extrémité ouverte.
Son coefficient de réflexion ne vaut que \Gneuf, et la tension en $m$ n'atteint que \rapNeuf{}
fois l'amplitude de la source, contre \rapMega{} avec \SI{e6}{\ohm}. Une partie importante de
l'onde incidente est absorbée par la charge à chaque arrivée.

% --- Amortissement ---
La ligne est sans pertes et seule la résistance de charge dissipe de l'énergie. Avec
\SI{e6}{\ohm} ou en court-circuit, $|\Gamma|$ est très proche de 1 et l'oscillation libre
devrait durer bien plus que 5~s. Elle disparaît pourtant en moins d'une seconde sur les
figures~\ref{fig:Q2_ouvert_5s} et~\ref{fig:Q3_cc_5s}. Cet amortissement est numérique : comme
$\tau/\Delta t$ n'est pas entier, la valeur retardée est interpolée linéairement, et cette
interpolation se comporte comme un filtre passe-bas à chaque passage de l'onde sur la ligne.
Choisir $\Delta t = \tau/N$ fait tomber le retard exactement sur un échantillon et supprime cet
effet. Une ligne réelle s'amortit, elle, par sa résistance linéique ; le cours propose d'en tenir
compte par une résistance localisée à chaque extrémité du modèle sans pertes.

% --- Court-circuit R = L = 0,001 ---
Le court-circuit $R = L = 0{,}001$ n'est pas parfait au pas de \SI{0.1}{ms}. Le modèle compagnon
de l'inductance vaut $R_{eq} = 2L/\Delta t = \ReqCC$~\unit{\ohm}, qui n'est pas négligeable devant
$Z_c$ : au premier pas, $v_2$ atteint encore \rapCC{} fois l'amplitude de la source. Ensuite $v_2$
change de signe à chaque pas. C'est l'oscillation numérique connue de la règle des trapèzes
quand le pas est grand devant la constante de temps du circuit ($L/Z_c$ est de quelques
microsecondes, $\Delta t$ vaut \SI{100}{\micro s}). Le courant, lui, n'est pas affecté de façon
visible.

% --- Source continue ---
Avec la source continue, le court-circuit ne présente pas de régime permanent fini : une source
de tension constante débite dans une résistance presque nulle à travers l'inductance de la ligne
et de la charge, et le courant croît donc sans cesse. Pour la charge R-L et pour 999~\unit{\ohm},
la réponse en escalier converge vers $V/R$. Avec la source sinusoïdale, $v_s$ varie pendant un
aller-retour, les marches se déforment et la réponse rejoint le régime donné par les phaseurs.

% --- Onde de choc ---
La formule de l'onde de choc n'atteint pas $V_{max}$ : sa crête est \chocCrete~kV, d'où la
correction $\eta = \etaVal$. Le front dure quelques microsecondes, bien moins que
$\tau = \tauL~\unit{\micro s}$ ; il faut donc un pas de quelques dizaines de nanosecondes, et la
fenêtre de \SI{250}{\micro s} est trop courte pour que l'onde traverse la ligne de 100~mi.
À cette échelle de temps, une inductance de charge se comporte presque en circuit ouvert, ce qui
explique le doublement de tension observé sur 10~mi avec la charge R-L.

\section{Conclusion}

Le modèle de Bergeron associé au modèle compagnon de la règle des trapèzes retrouve les valeurs
du livre ($i_2 = \ideuxC$~A à \SI{0.7}{ms}) et le régime permanent calculé par phaseurs. L'onde
atteint l'extrémité $m$ après $\tau = \tauL~\unit{\micro s}$. La tension y double presque à
l'extrémité ouverte (\rapMega{}\,$V_m$ avec \SI{e6}{\ohm}, mais seulement \rapNeuf{}\,$V_m$ avec
999~\unit{\ohm}) et s'effondre en court-circuit. Avec $\Delta t = \SI{0.1}{ms}$, l'interpolation
du retard amortit les oscillations libres ; un pas $\Delta t = \tau/N$ éviterait cet
amortissement artificiel. Pour le choc de foudre, la ligne de 100~mi n'a pas le temps de réagir
en \SI{250}{\micro s}, et il faut la raccourcir à 10~mi pour voir le doublement de tension
(\gVRL~kV).


\clearpage
\appendix
\section{Code MATLAB}

\lstinputlisting{TP01_ligne_longue.m}


\end{document}
